\documentclass[12pt, a4paper]{article}
\usepackage[utf8]{inputenc}
\usepackage[T1]{fontenc}
\usepackage{times}
\usepackage{geometry}
\usepackage{graphicx}
\usepackage{amsmath, amssymb, amsthm}
\usepackage{xcolor}
\usepackage{hyperref}
\usepackage{float}
\usepackage{titlesec}
\usepackage[bosnian]{babel}
\usepackage{fancyhdr}
\usepackage{chngcntr}
\usepackage{footnote}
\usepackage{threeparttable}

\makesavenoteenv{figure}  % allows normal footnotes inside figure captions

% Margins and formatting
\geometry{
 a4paper,
 left=25mm,
 right=25mm,
 top=30mm,
 bottom=25mm,
 headheight=15pt,
}

% Formatting specs
\setlength{\parskip}{1em}
\setlength{\parindent}{0pt}
\linespread{1.15}

% Figure numbering by section
\counterwithin{figure}{section}
\counterwithin{table}{section}
\counterwithin{equation}{section}

% Header and footer setup
\pagestyle{fancy}
\fancyhf{}
\fancyhead[L]{Numerička aproksimacija}
\fancyfoot[L]{Imamović, Jamaković, Habibović}
\fancyfoot[R]{\thepage}
\renewcommand{\headrulewidth}{0.4pt}
\renewcommand{\footrulewidth}{0.4pt}

% Bosnian captions
\renewcommand{\figurename}{Slika}
\renewcommand{\tablename}{Tabela}


\begin{document}

% ---------------- TITLE PAGE ----------------
\begin{titlepage}
    \thispagestyle{empty}
    \begin{center}
    \textbf{UNIVERZITET U ZENICI} \\
    \textbf{POLITEHNIČKI FAKULTET}
    \end{center}

    \vspace{6cm}

    \begin{center}
        \Large
        \textbf{NUMERIČKA APROKSIMACIJA} \\[0.3cm]
        \normalsize
        SEMINARSKI RAD
    \end{center}

    \vspace{8cm}

    \noindent
    \begin{minipage}[t]{0.45\textwidth}
        \textbf{STUDENTI} \\[0.3cm]
        Safet Imamović \\
        Emrah Jamaković \\
        Islam Habibović
    \end{minipage}
    \hfill
    \begin{minipage}[t]{0.45\textwidth}
        \raggedleft
        \textbf{PREDMETNI NASTAVNIK} \\[0.3cm]
        r. prof. dr. Aleksandar Karač
    \end{minipage}

    \vfill

    \begin{center}
        Zenica 2025/2026.
    \end{center}

\end{titlepage}

% ---------------- ABSTRACT ----------------
\pagenumbering{arabic}
\section*{Abstrakt}
U ovom radu detaljno je opisan proces razvoja i matematička pozadina softverskog rješenja za numeričku aproksimaciju funkcija. Fokus rada je na implementaciji \textit{Metode najmanjih kvadrata} (Least Squares Method) unutar modernog web okruženja. Aplikacija omogućava korisnicima unos eksperimentalnih podataka te pronalazak optimalne funkcije (linearne, kvadratne, polinomske, eksponencijalne ili stepene) koja najbolje opisuje trend podataka. Posebna pažnja posvećena je teorijskom izvođenju normalnih jednačina, analizi greške putem koeficijenta determinacije ($R^2$), te softverskoj arhitekturi baziranoj na Next.js i TypeScript tehnologijama. Rad sadrži i praktične primjere primjene razvijenog alata na realne inženjerske probleme.

\textbf{Ključne riječi:} Numerička aproksimacija, Metoda najmanjih kvadrata, Regresija, Next.js, TypeScript.

\newpage
\tableofcontents
\newpage

% ---------------- 1. UVOD ----------------
\section{UVOD}
U inženjerskoj i naučnoj praksi rijetko se susrećemo sa analitički zadanim funkcijama koje savršeno opisuju neki prirodni fenomen. Umjesto toga, na raspolaganju imamo skup diskretnih podataka dobijenih mjerenjem:
$$ (x_0, y_0), (x_1, y_1), \dots, (x_n, y_n) $$
Ovi podaci su često opterećeni mjernim šumom, greškama očitavanja ili nepravilnostima samog procesa. Zbog toga, pokušaj da se pronađe funkcija koja prolazi tačno kroz sve tačke (interpolacija) često daje nezadovoljavajuće rezultate, rezultirajući funkcijama koje naglo osciliraju između tačaka i nemaju fizikalnog smisla.

Rješenje ovog problema leži u \textbf{aproksimaciji}. Cilj aproksimacije nije pogoditi tačne vrijednosti u svakoj tački mjerenja, već pronaći "glatku" funkciju $f(x)$ koja prolazi \textit{što je moguće bliže} zadatim tačkama, minimizirajući ukupnu grešku. Najrasprostranjeniji metod za rješavanje ovog problema je \textbf{Metoda najmanjih kvadrata} (MNK), koju je početkom 19. vijeka neovisno razvio Carl Friedrich Gauss.

Ovaj seminarski rad ima dva osnovna cilja:
\begin{enumerate}
    \item Prikazati teorijsku osnovu metode najmanjih kvadrata, uključujući izvođenje formula za različite tipove funkcija.
    \item Opisati implementaciju ove metode u obliku moderne, interaktivne web aplikacije koja studentima i inženjerima služi kao alat za brzo modeliranje podataka.
\end{enumerate}

Kroz rad ćemo proći put od matematičkog formulisanja problema, preko rješavanja sistema linearnih jednačina u računarskom kodu, do vizualizacije rezultata u web pregledniku.

% ---------------- 2. TEORIJSKI DIO ----------------

\newpage

\section{TEORIJSKI DIO}

\subsection{Definicija problema aproksimacije}
Neka je dat skup od $N$ parova podataka $(x_i, y_i)$, gdje $i=1, \dots, N$. Tražimo funkciju $f(x; a_0, a_1, \dots, a_m)$ koja zavisi od parametara $a_k$, takvu da najbolje aproksimira zavisnost $y$ od $x$.
Definiramo \textbf{rezidual} (grešku) u svakoj tački kao razliku izmjerene i procijenjene vrijednosti:
\begin{equation}
    r_i = y_i - f(x_i)
\end{equation}

Osnovna ideja metode najmanjih kvadrata je minimizacija sume kvadrata ovih reziduala. Kvadrat se koristi kako bi se izbjeglo poništavanje pozitivnih i negativnih grešaka, te kako bi se veće greške "kaznile" više od malih. Funkcija kriterija $S$ je:
\begin{equation}
    S(a_0, \dots, a_m) = \sum_{i=1}^{N} (y_i - f(x_i))^2
\end{equation}
Optimalni parametri su oni za koje funkcija $S$ ima globalni minimum. Da bismo našli minimum, parcijalne izvode funkcije $S$ po svakom od parametara $a_k$ izjednačavamo sa nulom:
\begin{equation}
    \frac{\partial S}{\partial a_k} = 0, \quad k=0, \dots, m
\end{equation}

\subsection{Linearna regresija}
Najjednostavniji i najčešći oblik aproksimacije je linearna funkcija (pravac):
\begin{equation}
    y = ax + b
\end{equation}
Ovdje su parametri $a$ (nagib) i $b$ (odsječak). Funkcija greške je:
\begin{equation}
    S(a, b) = \sum_{i=1}^{N} (y_i - (ax_i + b))^2
\end{equation}
Minimizacija po $b$:
\begin{align}
    \frac{\partial S}{\partial b} &= \sum_{i=1}^{N} 2(y_i - ax_i - b)(-1) = 0 \\
    &= -2 \left( \sum y_i - a \sum x_i - \sum b \right) = 0 \\
    &\Rightarrow \sum y_i = a \sum x_i + N b
\end{align}
Minimizacija po $a$:
\begin{align}
    \frac{\partial S}{\partial a} &= \sum_{i=1}^{N} 2(y_i - ax_i - b)(-x_i) = 0 \\
    &= -2 \left( \sum x_i y_i - a \sum x_i^2 - b \sum x_i \right) = 0 \\
    &\Rightarrow \sum x_i y_i = a \sum x_i^2 + b \sum x_i
\end{align}
Dobili smo sistem od dvije linearne jednačine sa dvije nepoznate ($a$ i $b$), poznat kao \textbf{sistem normalnih jednačina}:
\begin{equation}
    \begin{bmatrix}
    N & \sum x_i \\
    \sum x_i & \sum x_i^2
    \end{bmatrix}
    \begin{bmatrix}
    b \\ a
    \end{bmatrix}
    =
    \begin{bmatrix}
    \sum y_i \\
    \sum x_i y_i
    \end{bmatrix}
\end{equation}

\subsection{Polinomska aproksimacija}
U općem slučaju, kada podatke želimo aproksimirati polinomom stepena $m$:
\begin{equation}
    P_m(x) = a_0 + a_1 x + a_2 x^2 + \dots + a_m x^m
\end{equation}
Postupak je analogan. Parcijalnim deriviranjem sume kvadrata grešaka dobija se sistem od $m+1$ linearnih jednačina. Matrica sistema $A$ sastoji se od suma stepena $x$:
\begin{equation}
    A = \begin{bmatrix}
    N & \sum x_i & \dots & \sum x_i^m \\
    \sum x_i & \sum x_i^2 & \dots & \sum x_i^{m+1} \\
    \vdots & \vdots & \ddots & \vdots \\
    \sum x_i^m & \sum x_i^{m+1} & \dots & \sum x_i^{2m}
    \end{bmatrix}
\end{equation}
Vektor slobodnih članova $B$ je:
\begin{equation}
    B = \begin{bmatrix}
    \sum y_i \\ \sum x_i y_i \\ \vdots \\ \sum x_i^m y_i
    \end{bmatrix}
\end{equation}
Rješavanjem sistema $A\mathbf{a} = B$ (npr. Gaussovom metodom) dobijamo koeficijente polinoma.

\subsection{Aproksimacija nelinearnim funkcijama (Linearizacija)}
Mnogi prirodni procesi (hlađenje, radioaktivni raspad, rast populacije) prate eksponencijalne ili stepene zakone, a ne polinomske. Primjena metode najmanjih kvadrata direktno na nelinearne funkcije dovela bi do sistema nelinearnih jednačina koji se teško rješavaju. Zato koristimo postupak \textbf{linearizacije}.

\subsubsection{Eksponencijalna funkcija}
Model: $y = a e^{bx}$.
Logaritmiranjem jednačine dobijamo:
\begin{equation}
    \ln y = \ln a + bx
\end{equation}
Uvodimo smjene varijabli:
$$ Y = \ln y, \quad A = \ln a, \quad B = b $$
Dobili smo linearnu jednačinu: $Y = A + Bx$. Sada na transformisanim podacima $(x_i, \ln y_i)$ primjenjujemo standardnu linearnu regresiju kako bismo našli $A$ i $B$. Konačni parametri su:
$$ a = e^A, \quad b = B $$

\subsubsection{Stepena funkcija}
Model: $y = a x^b$.
Logaritmiranjem:
\begin{equation}
    \ln y = \ln a + b \ln x
\end{equation}
Smjene:
$$ Y = \ln y, \quad X = \ln x, \quad A = \ln a, \quad B = b $$
Linearna forma: $Y = A + BX$. Parametri se određuju regresijom nad parovima $(\ln x_i, \ln y_i)$.

\subsection{Kvalitet aproksimacije}
Kako bismo ocijenili koliko dobro odabrana funkcija opisuje podatke, koristimo \textbf{Koeficijent determinacije} ($R^2$):
\begin{equation}
    R^2 = 1 - \frac{SS_{res}}{SS_{tot}} = 1 - \frac{\sum (y_i - f(x_i))^2}{\sum (y_i - \bar{y})^2}
\end{equation}
Gdje je $\bar{y}$ srednja vrijednost $y$ podataka.
\begin{itemize}
    \item $R^2 = 1$: Idealno poklapanje.
    \item $R^2 = 0$: Model ne objašnjava varijabilnost podataka.
\end{itemize}

% ---------------- 3. APLIKACIJA ----------------

\newpage

\section{APLIKACIJA}

U ovom poglavlju opisujemo tehničku realizaciju softverskog rješenja. Aplikacija je modularna platforma dizajnirana da bude proširiva i jednostavna za održavanje.

\subsection{Tehnološki okvir (Tech Stack)}
Prilikom izbora tehnologija vodili smo se kriterijima: brzina razvoja, performanse u pregledniku, i kvaliteta ekosistema.
\begin{itemize}
    \item \textbf{Next.js 15 (React Framework):} Omogućava nam kreiranje brzih, "server-side rendered" (SSR) ili statičkih web stranica. Iako se naša aplikacija u potpunosti izvršava na klijentu (SPA - Single Page Application), Next.js pruža odličnu strukturu projekta i routing.
    \item \textbf{TypeScript:} Superskup JavaScripta koji dodaje statičke tipove. U kontekstu numeričke matematike, ovo je neprocjenjivo. Na primjer, TypeScript nam ne dozvoljava da slučajno saberemo broj i niz, ili da funkciji koja očekuje matricu proslijedimo vektor.
    \item \textbf{MathLive:} Web komponenta (`<math-field>`) koja omogućava korisnicima unos formula na način kako se one pišu u matematici (WYSIWYG editor).
    \item \textbf{Cortex Compute Engine:} Moćan "engine" za parsiranje i evaluaciju simboličkih izraza. Koristimo ga da pretvorimo korisnikov unos (npr. "sin(x) + 2") u funkciju koja se može evaluirati za bilo koje $x$.
    \item \textbf{Plotly.js:} Industrijski standard za naučno iscrtavanje grafova na webu. Koristi WebGL za hardversku akceleraciju, što omogućava glatko iscrtavanje hiljada tačaka.
\end{itemize}

\subsection{Internacionalizacija (i18n)}

Aplikacija implementira internacionalizaciju (\textit{i18n}) kako bi bila
pristupačna korisnicima na više jezika. Trenutno su podržana dva jezika:
bosanski (\texttt{/bs}) i engleski (\texttt{/en}), pri čemu svaki jezik ima
svoju zasebnu rutu unutar aplikacije.

\subsection{Pokretanje aplikacije}

Za lokalno pokretanje aplikacije potrebno je imati instaliran
\textbf{Node.js} (verzija 18 ili novija) i \textbf{npm} (dolazi uz Node.js).
Nakon kloniranja repozitorija, postupak je sljedeći:

\begin{enumerate}
    \item Instalacija zavisnosti: \texttt{npm install}
    \item Pokretanje razvojnog servera: \texttt{npm run dev}
    \item Aplikacija je dostupna na \texttt{http://localhost:3000}
\end{enumerate}

Za produkcijski build koriste se komande \texttt{npm run build} i
\texttt{npm run start}. Aplikacija je također postavljena na
\href{https://aproksimacija.safet.dev}{\texttt{https://aproksimacija.safet.dev}}.

\subsection{Struktura izvornog koda}

Svi numerički algoritmi implementirani su kao čiste funkcije (\textit{pure functions})
u direktoriju \texttt{lib/math/}. Čiste funkcije za isti ulaz uvijek daju isti
izlaz i nemaju sporednih efekata, što ih čini pogodnim za verifikaciju i
testiranje. Ključni fajlovi su:

\begin{table}[H]
    \centering
    \begin{tabular}{|l|p{9cm}|}
        \hline
        \textbf{Fajl} & \textbf{Sadržaj} \\
        \hline
        \texttt{lib/types/index.ts} & Definicije tipova podataka (\texttt{DataPoint}, \texttt{ApproximationResult}, \texttt{InterpolationResult}, \texttt{MetodaRjesavanja}) \\
        \hline
        \texttt{lib/math/aproksimacija.ts} & Pet metoda aproksimacije: linearna, kvadratna, polinomna, stepena i eksponencijalna \\
        \hline
        \texttt{lib/math/interpolation.ts} & Tri metode interpolacije: Lagrangeova, Newtonova i direktna \\
        \hline
        \texttt{lib/math/matrica-utils.ts} & Operacije linearne algebre: Gaussova eliminacija, Gauss-Jordan (RREF), LU faktorizacija (Doolittle), te formatiranje u LaTeX \\
        \hline
        \texttt{lib/math/expression-parser.ts} & Parsiranje i evaluacija matematičkih izraza (omotač oko Cortex Compute Engine) \\
        \hline
        \texttt{lib/math/validators.ts} & Validacija ulaznih podataka \\
        \hline
    \end{tabular}
    \caption{Ključni fajlovi sa numeričkim algoritmima}
\end{table}

\subsection{Rješavanje sistema linearnih jednačina}

Sve metode aproksimacije i direktna interpolacija svode se na
rješavanje sistema oblika $A\mathbf{x} = \mathbf{b}$. Aplikacija
nudi korisniku izbor između tri metode rješavanja, implementirane
u fajlu \texttt{lib/math/matrica-utils.ts}:

\begin{enumerate}
    \item \textbf{Gaussova eliminacija} --- parcijalno pivotiranje,
    direktna eliminacija do gornje-trougaonog oblika, povratna supstitucija.
    \item \textbf{Gauss-Jordan eliminacija} --- proširuje Gaussovu metodu
    normalizacijom pivot reda i eliminacijom iznad pivota, čime se dobija
    reducirani redni ešalonski oblik (RREF). Rješenje se čita direktno
    iz posljednje kolone.
    \item \textbf{LU faktorizacija (Doolittle)} --- dekompozicija matrice
    $A = L \cdot U$ na donje-trougaonu ($L$) i gornje-trougaonu ($U$)
    matricu, nakon čega se rješavaju dva trougaona sistema:
    $L\mathbf{y'} = \mathbf{b}$ i $U\mathbf{x} = \mathbf{y'}$.
\end{enumerate}

Svaka metoda generira detaljan prikaz koraka rješavanja u LaTeX formatu,
koji se prikazuje korisniku u sekciji ``Koraci izračuna''.

\subsection{Implementacija aproksimacije}

Metode aproksimacije nalaze se u fajlu \texttt{lib/math/aproksimacija.ts}.
Svaka metoda izračunava potrebne sume ($\sum x_i^k$, $\sum x_i^k y_i$),
formira sistem normalnih jednačina, te ga prosljeđuje odabranom solveru.

Za \textbf{polinomnu aproksimaciju} stepena $m$, normalne jednačine se
formiraju direktno iz suma, u skladu sa formulama iz poglavlja 2.3:
\[
    \begin{bmatrix}
    N & \sum x_i & \cdots & \sum x_i^m \\
    \sum x_i & \sum x_i^2 & \cdots & \sum x_i^{m+1} \\
    \vdots & \vdots & \ddots & \vdots \\
    \sum x_i^m & \sum x_i^{m+1} & \cdots & \sum x_i^{2m}
    \end{bmatrix}
    \begin{bmatrix} a_0 \\ a_1 \\ \vdots \\ a_m \end{bmatrix}
    =
    \begin{bmatrix}
    \sum Y_i \\ \sum x_i Y_i \\ \vdots \\ \sum x_i^m Y_i
    \end{bmatrix}
\]

Za \textbf{nelinearne modele} (stepena i eksponencijalna funkcija)
koristi se linearizacija logaritmovanjem, nakon čega se problem
svodi na linearnu regresiju nad transformisanim podacima.

\subsection{Implementacija interpolacije}

Metode interpolacije nalaze se u fajlu \texttt{lib/math/interpolation.ts}:

\begin{itemize}
    \item \textbf{Lagrangeova interpolacija} --- izračunava bazne polinome
    $L_k(x)$ i razvija ih u standardnu formu.
    \item \textbf{Newtonova interpolacija} --- gradi tablicu podijeljenih
    razlika i konstruira polinom u Newtonovoj formi, koji se potom
    pretvara u standardni oblik.
    \item \textbf{Direktna interpolacija} --- formira sistem jednačina
    koristeći matricu koeficijenata i rješava ga odabranom metodom
    (Gauss, Gauss-Jordan ili LU).
\end{itemize}

Sve tri metode za iste ulazne podatke daju isti interpolacijski polinom.

\subsection{Validacija ulaza i numerička stabilnost}

Jedan od ključnih izazova u implementaciji numeričkih metoda u softveru
jeste osiguravanje numeričke stabilnosti i sprječavanje neispravnih
kombinacija ulaznih podataka. Iako matematički modeli mogu biti jasno
definisani, korisnik u praktičnoj aplikaciji može unijeti podatke koji
ne zadovoljavaju teorijske pretpostavke metode.

Posebna pažnja posvećena je aproksimaciji polinomima višeg stepena.
Polinom stepena $m$ ima $m+1$ nepoznatih koeficijenata, te je za rješavanje
sistema normalnih jednačina neophodno imati najmanje $m+1$ nezavisnih
tačaka. U suprotnom, sistem postaje singularan i nema jedinstveno rješenje.

Zbog toga aplikacija automatski ograničava maksimalni stepen polinoma
na osnovu broja unesenih tačaka:
\[
m_{\text{max}} = N - 1
\]
gdje je $N$ broj dostupnih podataka. Ova provjera implementirana je
na dva nivoa:
\begin{itemize}
    \item \textbf{Korisničko sučelje (UI):} Stepen polinoma se automatski
    prilagođava prilikom dodavanja ili uklanjanja tačaka, čime se korisniku
    onemogućava izbor nevalidnih konfiguracija.
    \item \textbf{Logika izračuna:} Prije samog proračuna provjerava se da li
    postoji dovoljan broj tačaka za izabranu metodu, čime se sprječava
    numerička nestabilnost i pad aplikacije.
\end{itemize}

Slične validacije primijenjene su i kod kvadratne aproksimacije, gdje je
minimalan broj tačaka tri, kao i kod nelinearnih modela (eksponencijalni i
stepeni), gdje se dodatno provjerava domen funkcija (pozitivne vrijednosti
za logaritamsku transformaciju).

\subsection{Načini unosa podataka}

Aplikacija podržava tri načina unosa podatkovnih tačaka:

\begin{enumerate}
    \item \textbf{Ručni unos} --- korisnik direktno upisuje parove $(x_i, y_i)$
    u tabelu. Redovi se mogu dodavati i uklanjati po potrebi. Ovo je
    najjednostavniji način za manji broj tačaka.

    \item \textbf{Učitavanje fajla} --- korisnik može učitati podatke iz
    eksternog fajla metodom \textit{drag-and-drop} ili odabirom fajla.
    Podržani formati su:
    \begin{itemize}
        \item \textbf{CSV} --- vrijednosti razdvojene zarezom (\texttt{x,y})
        ili tačka-zarezom (\texttt{x;y}),
        \item \textbf{DAT} --- vrijednosti razdvojene tabulatorom (\texttt{x\textbackslash ty}),
        \item \textbf{JSON} --- niz objekata oblika \texttt{[\{x, y\}, ...]}.
    \end{itemize}

    \item \textbf{Funkcijski unos} --- korisnik unosi matematički izraz
    $f(x)$ putem \textit{MathLive} editora, koji omogućava LaTeX-stilski
    unos formula (npr.\ $\sin(x) + x^2$). Nakon definisanja domene
    $[x_{\min}, x_{\max}]$ i broja uzoraka, aplikacija automatski generira
    tačke evaluacijom izraza. Ovaj način je posebno koristan za testiranje
    metoda na poznatim funkcijama.
\end{enumerate}

Pored toga, dostupni su i unaprijed pripremljeni primjeri skupova podataka
iz različitih inženjerskih oblasti (Hookeov zakon, pražnjenje kondenzatora,
karakteristika pumpe i dr.), koje korisnik može učitati jednim klikom.

\newpage

\subsection{Korisničko sučelje (UI)}

Aplikacija je dostupna na
\href{https://aproksimacija.safet.dev/bs}{\texttt{https://aproksimacija.safet.dev/bs}},
ili lokalno na \texttt{http://localhost:3000/bs} nakon pokretanja
razvojnog servera. Implementirane su metode aproksimacije i interpolacije,
sa planom proširenja na diferenciranje i integriranje.

\begin{figure}[H]
    \centering
    \includegraphics[width=1\textwidth]{landing_page.png}
    \caption{Početna stranica aplikacije za Numeričku Aproksimaciju.}
    \label{fig:ui-polynomial}
\end{figure}


Aplikacija koristi tzv. "Split-screen" raspored. Gornja polovica ekrana rezervirana je za unos podataka (tabela) i kontrolu parametara, dok donja polovica prikazuje rezultate (graf i tekstualno objašnjenje).
Komponente su stilizirane koristeći `Tailwind CSS` klase, što nam omogućava da bez pisanja posebnih `.css` fajlova definiramo izgled elemenata.

\begin{figure}[H]
    \centering
    \includegraphics[width=0.8\textwidth]{Sucelje.png}
    \caption{Korisničko sučelje aplikacije za polinomsku aproksimaciju}
    \label{fig:ui-polynomial}
\end{figure}

% ---------------- 4. PRIMJERI PRIMJENE ----------------
\newpage

\section{PRIMJERI PRIMJENE}

U ovom poglavlju testiramo aplikaciju na nekoliko različitih scenarija, koji pokrivaju različite vrste aproksimacije.

Postoje primjeri skupova podataka koje korisnik može odabrati.

\begin{figure}[H]
    \centering
    \includegraphics[width=1\textwidth]{primjeri_podaci.png}
    \caption{Primjeri Podataka.}
    \label{fig:ui-polynomial}
\end{figure}

\subsection{Primjer 1: Hookeov zakon (Fizika)}
Testiramo linearnu aproksimaciju na primjeru opruge. Mjerimo produženje opruge ($\Delta x$) pod djelovanjem različitih sila ($F$). Teoretski zakon je linearan: $F = k \cdot \Delta x$.

\textbf{Ulazni podaci:}
\begin{table}[H]
    \centering
    \begin{tabular}{|c|c|c|c|c|c|}
        \hline
        Produženje $\Delta x$ [cm] & 2 & 4 & 6 & 8 & 10 \\
        \hline
        Sila $F$ [N] & 5.1 & 9.8 & 15.2 & 19.9 & 25.1 \\
        \hline
    \end{tabular}
\end{table}

\begin{figure}[H]
    \centering
    \includegraphics[width=0.9\textwidth]{Unos1.png}
    \caption{Prikaz unesenih podataka u aplikaciju}
\end{figure}





Unošenjem ovih podataka u našu aplikaciju i odabirom \textit{Linear Approximation}, dobijamo sljedeći rezultat:
\begin{itemize}
    \item \textbf{Funkcija:} $F(x) = 2.505x + 0.01$
    \item \textbf{Koeficijent krutosti:} $k \approx 2.505$ N/cm.
    \item \textbf{Kvalitet:} $R^2 = 0.9996$, što ukazuje na izuzetno slaganje sa linearnim modelom.
\end{itemize}

\begin{figure}[H]
    \centering
    \includegraphics[width=0.9\textwidth]{Primjer1.png}
    \caption{Grafički prikaz linearne aproksimacije za Hookeov zakon}
\end{figure}

\begin{figure}[H]
    \centering
    \includegraphics[width=0.8\textwidth]{lin_koraci.png}
    \caption{Grafički prikaz rezultata i koraka linearne aproksimacije za Hookeov zakon}
    \label{fig:exp-capacitor}
\end{figure}

\begin{figure}[H]
    \centering
    \includegraphics[width=0.9\textwidth]{line_reg.png}
    \caption{Validacija linearne aproksimacije za Hookeov zakon putem stats blue Stats Suite\textsuperscript{\cite{statsblue}}.}
\end{figure}



\subsection{Primjer 2: Kondenzator (Elektrotehnika)}
Posmatramo pražnjenje kondenzatora kroz otpornik. Napon opada po eksponencijalnom zakonu $V(t) = V_0 e^{-t/RC}$.

\textbf{Ulazni podaci (Vrijeme $t$ [s] vs Napon $V$ [V]):}
\begin{table}[H]
    \centering
    \begin{tabular}{|c|c|c|c|c|c|}
        \hline
        $t$ [s] & 0 & 1 & 2 & 3 & 4 \\
        \hline
        $V$ [V] & 10.0 & 6.1 & 3.7 & 2.2 & 1.4 \\
        \hline
    \end{tabular}
\end{table}

\begin{figure}[H]
    \centering
    \includegraphics[width=0.9\textwidth]{Unos2.png}
    \caption{Prikaz unesenih podataka u aplikaciju}
\end{figure}

Odabirom \textit{Exponential Approximation} u aplikaciji:
\begin{itemize}
    \item \textbf{Transformacija:} Aplikacija interno logaritmira napone ($\ln V$) i vrši linearnu regresiju nad vremenom.
    \item \textbf{Rezultat:} $V(t) \approx 9.9665918088 \cdot e^{−0.4952057124t}$
    \item \textbf{Analiza:} Početni napon je približno 10V (što odgovara podacima), a vremenska konstanta $\tau = 1/0.498 \approx 2.0$s.
\end{itemize}

\begin{figure}[H]
    \centering
    \includegraphics[width=0.7\textwidth]{Primjer2.png}
    \caption{Grafički prikaz eksponencijalne aproksimacije pražnjenja kondenzatora}
    \label{fig:exp-capacitor}
\end{figure}

\begin{figure}[H]
    \centering
    \includegraphics[width=0.8\textwidth]{exp_koraci.png}
    \caption{Grafički prikaz rezultata i koraka rješenja eksponencijalne aproksimacije pražnjenja kondenzatora}
    \label{fig:exp-capacitor}
\end{figure}

\begin{figure}[H]
    \centering
    \includegraphics[width=0.6\textwidth]{exp.png}
    \caption{Validacija eksponencijalne aproksimacije pražnjenja kondenzatora putem stats blue Stats Suite\textsuperscript{\cite{statsblue}}.}
\end{figure}

\newpage

\subsection{Primjer 3: Karakteristika pumpe (Mašinstvo)}
Karakteristika centrifugalne pumpe (zavisnost visine dizanja $H$ od protoka $Q$) je obično kvadratna funkcija: $H(Q) = H_0 - kQ^2$.

\textbf{Ulazni podaci:}
\begin{table}[H]
    \centering
    \begin{tabular}{|c|c|c|c|c|c|}
        \hline
        Protok $Q$ [$m^3/h$] & 0 & 10 & 20 & 30 & 40 \\
        \hline
        Visina $H$ [m] & 50.0 & 49.2 & 46.5 & 42.1 & 35.5 \\
        \hline
    \end{tabular}
\end{table}

\begin{figure}[H]
    \centering
    \includegraphics[width=0.9\textwidth]{Unos3.png}
    \caption{Prikaz unesenih podataka u aplikaciju}
\end{figure}



Korištenjem \textit{Quadratic Approximation} (polinom 2. reda):
\begin{itemize}
    \item \textbf{Rezultat:} $H(Q) = 49.98 + 0.019 Q - 0.0095 Q^2$
    \item \textbf{Greška:} $R^2 = 0.999$, model odlično prati eksperimentalne podatke.
\end{itemize}
Korisnik sada može pomoću aplikacije predvidjeti visinu dizanja za protok od npr. $Q=25 m^3/h$ jednostavnim pozicioniranjem kursora na graf.

\begin{figure}[H]
    \centering
    \includegraphics[width=0.8\textwidth]{Primjer3.png}
    \caption{Grafički prikaz kvadratne aproksimacije karakteristike pumpe}
    \label{fig:pump-quadratic}
\end{figure}

\begin{figure}[H]
    \centering
    \includegraphics[width=0.8\textwidth]{Result3.png}
    \caption{Grafički prikaz rezultata i koraka rješenja kvadratne aproksimacije karakteristike pumpe}
\end{figure}

\begin{figure}[H]
    \centering
    \includegraphics[width=0.9\textwidth]{3_red_2.png}
    \caption{Validacija kvadratne aproksimacije karakteristike pumpe putem stats blue Stats Suite\textsuperscript{\cite{statsblue}}.}
\end{figure}

\newpage

\subsection{Primjer 4: Balistička putanja projektila (Fizika)}

Razmatramo horizontalno ispaljen projektil, gdje se mjeri zavisnost
visine projektila $y$ od horizontalne udaljenosti $x$.
Zbog djelovanja gravitacije, putanja projektila ima paraboličan oblik.

\textbf{Ulazni podaci:}

\begin{table}[H]
    \centering
    \begin{tabular}{|c|c|c|c|c|c|}
        \hline
        Horizontalna udaljenost $x$ [m] & 0 & 10 & 20 & 30 & 40 \\
        \hline
        Visina $y$ [m] & 5.0 & 7.8 & 8.9 & 8.2 & 5.7 \\
        \hline
    \end{tabular}
\end{table}

\begin{figure}[H]
    \centering
    \includegraphics[width=0.9\textwidth]{Unos4.png}
    \caption{Prikaz unesenih podataka u aplikaciju}
\end{figure}

Korištenjem \textit{Polynomial Approximation} u aplikaciji (polinom
drugog stepena) dobijamo aproksimacioni model:

\begin{itemize}
    \item \textbf{Rezultat:}
    \[
    y(x) \approx 4.9885714286 + 0.3722857143x - 0.0088571429x^2
    \]
    \item \textbf{Analiza:} Negativan koeficijent uz $x^2$ odgovara
    fizičkom očekivanju opadanja visine pod uticajem gravitacije.
    \item \textbf{Greška:} $R^2 = 0.997$, što ukazuje na vrlo dobro
    slaganje modela sa eksperimentalnim podacima.
\end{itemize}

\begin{figure}[H]
    \centering
    \includegraphics[width=0.8\textwidth]{Primjer4.png}
    \caption{Grafički prikaz polinomne aproksimacije putanje projektila}
    \label{fig:pump-quadratic}
\end{figure}

\begin{figure}[H]
    \centering
    \includegraphics[width=0.8\textwidth]{Result4_.png}
    \caption{Grafički prikaz rezultata i koraka rješenja polinomne aproksimacije putanje projektila}
    \label{fig:pump-quadratic}
\end{figure}

\begin{figure}[H]
    \centering
    \includegraphics[width=0.8\textwidth]{polinom.png}
    \caption{Validacija polinomne aproksimacije putanje projektila putem stats blue Stats Suite\textsuperscript{\cite{statsblue}}.}
\end{figure}

Polinomska aproksimacija omogućava jednostavno modeliranje putanje
projektila i predviđanje maksimalne visine ili dometa.

\subsection{Primjer 5: Otpor fluida u zavisnosti od brzine (Mehanika fluida)}

Kod kretanja tijela kroz fluid (zrak ili tečnost), sila otpora
zavisi od brzine kretanja. U mnogim praktičnim slučajevima,
ova zavisnost se opisuje stepenim zakonom:
\[
F = a v^b
\]

\textbf{Ulazni podaci:}

\begin{table}[H]
    \centering
    \begin{tabular}{|c|c|c|c|c|c|}
        \hline
        Brzina $v$ [m/s] & 1 & 2 & 3 & 4 & 5 \\
        \hline
        Sila otpora $F$ [N] & 0.5 & 2.1 & 4.8 & 8.9 & 14.2 \\
        \hline
    \end{tabular}
\end{table}

\begin{figure}[H]
    \centering
    \includegraphics[width=0.9\textwidth]{Unos5.png}
    \caption{Prikaz unesenih podataka u aplikaciju}
\end{figure}

U aplikaciji se bira \textit{Power Approximation}, pri čemu se vrši
logaritamska transformacija podataka i primjenjuje metoda najmanjih
kvadrata.

\begin{itemize}
    \item \textbf{Rezultat:}
    \[
    F(v) \approx 0.4976907413 \cdot v^{2.0774095997}
    \]
    \item \textbf{Analiza:} Eksponent $b \approx 1.92$ blizak je
    kvadratnoj zavisnosti, što odgovara teoriji otpora fluida
    pri većim brzinama.
    \item \textbf{Greška:} $R^2 = 0.998$, što potvrđuje vrlo dobar
    kvalitet aproksimacije.
\end{itemize}

\begin{figure}[H]
    \centering
    \includegraphics[width=0.8\textwidth]{Primjer5.png}
    \caption{Grafički prikaz stepene aproksimacije otpora fluida}
    \label{fig:pump-quadratic}
\end{figure}

\begin{figure}[H]
    \centering
    \includegraphics[width=0.65\textwidth]{Result5.png}
    \caption{Grafički prikaz rezultata i koraka rješenja stepene aproksimacije otpora fluida}
    \label{fig:pump-quadratic}
\end{figure}

\begin{figure}[H]
    \centering
    \includegraphics[width=0.32\textwidth]{power_koraci.png}
    \caption{Validacija kvadratne aproksimacije karakteristike pumpe putem Statology\textsuperscript{\cite{statology}}.}
\end{figure}

Stepena aproksimacija omogućava jednostavno modeliranje sile otpora
i predviđanje ponašanja sistema pri različitim brzinama kretanja.
\newpage

% ---------------- 5. ZAKLJUČAK ----------------
\section{ZAKLJUČCI}
Cilj ovog seminarskog rada bio je dvojak: teoretski obraditi problem numeričke aproksimacije metodom najmanjih kvadrata i praktično implementirati te metode u modernom softverskom okruženju.

Kroz razvoj aplikacije pokazali smo da:
\begin{enumerate}
    \item \textbf{Linearna algebra je temelj:} Implementacija robusnih rješavača linearnih sistema (Gaussova eliminacija, Gauss-Jordan, LU faktorizacija) je ključna za sve metode aproksimacije i interpolacije.
    \item \textbf{Linearizacija je moćan alat:} Svođenje nelinearnih problema (eksponencijalni, stepeni) na linearne putem logaritamske transformacije drastično pojednostavljuje proračun i omogućava korištenje iste baze koda.
    \item \textbf{Web tehnologije su zrele za inženjerske alate:} Sposobnost JavaScript-a (i TypeScript-a) da brzo manipulira nizovima i matricama, u kombinaciji sa bibliotekama za iscrtavanje poput Plotly.js, čini ga odličnim izborom za razvoj inženjerskog softvera.
\end{enumerate}

Aplikacija uspješno ispunjava sve postavljene zahtjeve: tačna je, brza, i pruža vizualni uvid u kvalitet aproksimacije i interpolacije. Korisnik može birati metodu rješavanja linearnog sistema, pratiti svaki korak postupka u LaTeX formatu, te pregledavati historiju izračuna. Kao takva, predstavlja koristan edukativni alat za studente Politehničkog fakulteta.

\newpage

\begin{thebibliography}{9}  % the number here is the widest label
\bibitem{chapra2015} S. C. Chapra, R. P. Canale, \textit{Numerical Methods for Engineers}, 7th Edition, McGraw-Hill Education, 2015.
\bibitem{karac2026} A. Karač, \textit{Predavanja iz predmeta Primijenjene numeričke metode u softverskom inženjerstvu}, Politehnički fakultet, Univerzitet u Zenici, 2026.
\bibitem{vercel} Vercel Inc., \textit{Next.js Documentation}, \url{https://nextjs.org}, pristupljeno: Februar 2026.
\bibitem{plotly} Plotly Technologies Inc., \textit{Plotly JavaScript Open Source Graphing Library}, \url{https://plotly.com/javascript/}, pristupljeno: Februar 2026.
\bibitem{mathlive} MathLive, \textit{A Web Component for Math Input}, \url{https://cortexjs.io/mathlive/}, pristupljeno: Februar 2026.
\bibitem{statsblue} Stats Blue, \textit{Stats Suite}, \url{https://stats.blue/}, pristupljeno: Februar 2026.
\bibitem{statology} Zach Bobbitt, \textit{Power Regression Calculator}, \url{https://www.statology.org/power-regression-calculator/}, pristupljeno: Februar 2026.

\end{thebibliography}

\end{document}
